\chapter{Quantum sets and quantum graphs}

In this chapter we will introduce the notions of quantum sets and quantum
graphs. They were first considered by...

\section{Quantum sets}

Let \( B \) be a \(C^*\)-algebra equipped with a faithful positive functional
\( \psi \). By the Gelfand-Naimark theorem, \( B \) can be realized concretely
as a closed \(*\)-subalgebra of the bounded linear operators \( B(\mathcal{H})
\) on some Hilbert space \( \mathcal{H} \). When \( B \) is finite-dimensional,
this abstract framework simplifies considerably: we can identify \( B \) with
its own Hilbert space completion \( L^2(B, \psi) \) by defining the inner
product directly via the functional: \[ \langle x, y \rangle = \psi(x^* y)
\quad \forall x, y \in B. \]

The following definition introduces one of the central notions used
in this thesis. It is taken directly from \textcite{tobolski26}.

\begin{definition}[Quantum set, \textcite[Definition~2.1]{tobolski26}]
    A quantum set \( (B, \psi) \) is a pair consisting of a finite-dimensional
    \(C^{*}\)-algebra \( B \) and a faithful positive functional \( \psi\colon B
    \rightarrow \mathbb{C} \) such that
    \[
        mm^{*} = \operatorname{Id}_{B}
    \]
    holds for the multiplication map \( m\colon B \otimes B \rightarrow B
    \) and its adjoint \( m^{*} \) given by
    \[
        \left\langle m(a \otimes b), c \right\rangle = \left\langle a
        \otimes b, m^{*}(c) \right\rangle_{\psi \otimes \psi}
    \]
    for all \( a, b, c \in B \).
\end{definition}

\begin{remark}
    We will call any functional with the property given in the
    definition above a \emph{quantum set functional} on \( B \).
\end{remark}

\begin{remark}
    Recall that in general, a \emph{comultiplication} on a
    \(C^{*}\)-algebra is a \(*\)-homomorphism
    \[
        \Delta\colonA\to A\otimes A
    \]
    such that the following diagram commutes:
    \[
        \begin{tikzcd}
            A \arrow[r,"\Delta"] \arrow[d,"\Delta"']
            & A\otimes A \arrow[d,"{\Delta\otimes\mathrm{id}}"] \\
            A\otimes A \arrow[r,"{\mathrm{id}\otimes\Delta}"']
            & A\otimes A\otimes A
        \end{tikzcd}
    \]
    Note that in general a comultiplication is \emph{not} unique. In our
    setting it is uniquely determined by the additional structure and the
    Riesz representation theorem for Hilbert spaces.
\end{remark}

\begin{remark}
    Since in our setup the comultiplication is given by the scalar product
    and regular multiplication \( m \), we will denote it by \( m^{*} \).
\end{remark}

% TODO: Prove that the examples are correct (do some calculations)
\begin{example}
    Every finite set \( X \) gives rise to a quantum set \( (C(X), \psi_{X}) \)
    with \( \psi_{X}\colon C(X) \rightarrow \mathbb{C} \), \(
        \psi_{X}(f) = \sum_{x
    \in X} f(x) \) for any \( f \in C(X) \).
    Let \( X \) be a finite set and \( C(X) \) be the \(*\)-algebra of
    complex-valued functions on \( X \), where the involution is given by
    pointwise complex conjugation: \( f^*(x) = \overline{f(x)} \).

    To show that \( \psi_X \) is positive, we must verify that
    \( \psi_X(f^* f) \geq 0 \) for all \( f \in C(X) \). For any function
    \( f \), the product \( f^* f \) evaluates pointwise to:
    \[
        (f^* f)(x) = \overline{f(x)} f(x) = |f(x)|^2.
    \]
    Applying the functional yields:
    \[
        \psi_X(f^* f) = \sum_{x \in X} |f(x)|^2.
    \]
    Thus the functional is positive.

    To prove faithfulness, suppose \( \psi_X(f^* f) = 0 \) for some \( f
    \in C(X) \). This implies:
    \[
        \sum_{x \in X} |f(x)|^2 = 0.
    \]
    Because each term \( |f(x)|^2 \) is strictly non-negative, a sum of
    such terms can only equal zero if every individual term is zero.
    Thus, we must have:
    \[
        |f(x)|^2 = 0 \quad \forall x \in X,
    \]
    which directly implies \( f(x) = 0 \) for all \( x \in X \). Therefore,
    \( f \) is the zero function (\( f = 0 \)), establishing that \(
    \psi_X \) is faithful.

    Let \( \{e_x\}_{x \in X} \) be the standard orthonormal basis of
    characteristic functions on \( X \). Recall the definitions of the
    multiplication map \( m \) and its adjoint \( m^* \):
    \begin{align*}
        m(f \otimes g) &= fg, \\
        m^*(e_x) &= e_x \otimes e_x.
    \end{align*}

    We wish to show that the composition \( mm^*\colon C(X)
    \rightarrow C(X) \) is the identity operator (\( \text{Id} \)). Let us
    evaluate the action of \( m m^* \) on an arbitrary basis element \( e_x \):
    \[
        (m m^*)(e_x) = m\big(m^*(e_x)\big) = m(e_x \otimes e_x) = e_x e_x = e_x.
    \]

    Because the linear operators \( m m^* \) and \( \text{Id} \) agree on
    every element of the basis \( \{e_x\}_{x \in X} \), by linearity they
    must be equal on the entire space \( C(X) \). Therefore:
    \[
        m m^* = \text{Id}.
    \]
\end{example}

% TODO: Is this necessary?
\begin{example}
    In particular, \( (\mathbb{C}^{n}, \psi_{n}) \) with \( \psi_{n}(e_{i}) = 1
    \) is a quantum set. The functional naturally extends to:
    \[
        \psi_n(x) = \sum_{i=1}^n x_i.
    \]

    To show that \(\psi_n\) is positive, we evaluate it on the
    element \(x^* x\). Using the orthogonality of the basis, we have:
    \[
        x^* x = \left( \sum_{i=1}^n \overline{x_i} e_i \right) \left(
        \sum_{j=1}^n x_j e_j \right) = \sum_{i,j=1}^n \overline{x_i}
        x_j e_i e_j = \sum_{i=1}^n \overline{x_i} x_i e_i =
        \sum_{i=1}^n |x_i|^2 e_i.
    \]
    Applying the functional yields:
    \[
        \psi_n(x^* x) = \psi_n\left( \sum_{i=1}^n |x_i|^2 e_i \right)
        = \sum_{i=1}^n |x_i|^2 \psi_n(e_i) = \sum_{i=1}^n |x_i|^2.
    \]
    Since \(|x_i|^2 \geq 0\) for all \(i\), the sum of these
    non-negative real numbers is also non-negative, meaning
    \(\psi_n(x^* x) \geq 0\). Thus, the functional is positive.

    To prove faithfulness, suppose \(\psi_n(x^* x) = 0\) for some \(x
    \in \mathbb{C}^n\). This implies:
    \[
        \sum_{i=1}^n |x_i|^2 = 0.
    \]
    Because each term \(|x_i|^2\) is strictly non-negative, the sum
    can only be zero if every individual component satisfies
    \(|x_i|^2 = 0\) for all \(i \in \{1, \dots, n\}\). This directly
    implies \(x_i = 0\) for all \(i\). Therefore, \(x = \sum_{i=1}^n
    0 \cdot e_i = 0\), establishing that \(\psi_n\) is faithful.

    The adjoint map \(m^*\colon \mathbb{C}^n
    \rightarrow \mathbb{C}^n \otimes \mathbb{C}^n\) acts on a basis
    element \(e_i\) by:
    \[
        m^*(e_i) = e_i \otimes e_i.
    \]

    We wish to show that the composition \(mm^*\colon \mathbb{C}^n
    \rightarrow \mathbb{C}^n\) is the identity operator. We check:
    \[
        (m m^*)(e_i) = m\big(m^*(e_i)\big) = m(e_i \otimes e_i) = e_i.
    \]
    Because the linear operators \(m m^* \) and \(\text{Id}\) agree
    on every element of the basis \(\{e_i\}_{i=1}^n\), by linearity
    they must be equal on the entire vector space \(\mathbb{C}^n\). Therefore:
    \[
        m m^* = \text{Id}.
    \]
\end{example}

% TODO: Why is there nTr instead of Tr?
\begin{example}
    For any \( n \in \mathbb{N}_{+} \), \( (M_{n}, n
    \operatorname{Tr}) \) is a quantum set, with \( \operatorname{Tr}
    \) being the trace. We will later show that this structure can be
    obtained as the crossed product \( (\mathbb{C}^{n}, \psi_{n}) \) and
    \( \hat{\mathbb{Z}}_{n} \) with the standard action, i.e.
    \[
        (M_{n}, n \operatorname{Tr}) \cong (\mathbb{C}^{n} \rtimes
        \hat{\mathbb{Z}}_{n}, \psi_{\rtimes}).
    \]
\end{example}

% TODO: give a better name
\begin{lemma}[Comultiplication for matrix algebras]
    In the context of \( (M_{n}, n \operatorname{Tr}) \), the
    comultiplication
    is given by the formula
    \[
        m^{*}(e_{ij}) = \sum_{k=1}^{n} e_{ik} \otimes e_{kj}
    \]
    with \( e_{ij} \) being the matrix with a single \( 1 \) in
    the \( i \)-th
    row and \( j \)-th column, and extended linearly.
\end{lemma}

% TODO: check this proof carefully
\begin{proof}
    This is simply a computational excercise. We will use the following
    well-known formulas:
    \begin{enumerate}
        \item \( e_{ab}e_{cd} = \delta_{bc}e_{ad} \),
        \item \( \left\langle e_{ij}, e_{kl} \right\rangle =
            \operatorname{Tr}(e_{ji}e_{kl}) = \delta_{ik}e_{jl} \),
    \end{enumerate}
    with
    \[
        \delta_{ij} =
        \begin{cases}
            1 & \text{if } i = j,\\
            0 & \text{otherwise}.
        \end{cases}
    \]
    In particular \( \left\{ e_{ij} \right\}_{i,j=1,\dots,n} \) form
    a orthonormal basis.

    Take arbitrary basis vectors \( e_{ab} \), \( e_{cd} \). Then
    \[
        m(e_{ab} \otimes e_{cd}) = e_{ab} e_{cd} = \delta_{bc}e_{ad},
    \]
    hence
    \begin{align*}
        \left\langle m(e_{ab} \otimes e_{cd}), e_{ij} \right\rangle
        &= \delta_{bc} \left\langle e_{ad}, e_{ij} \right\rangle \\
        &= \delta_{bc}\delta_{ai}\delta_{dj}.
    \end{align*}
    Now compute the inner product with the proposed adjoint:
    \begin{align*}
        \left\langle e_{ab} \otimes e_{cd}, \sum_{k=1}^{n} e_{ik}
        \otimes e_{kj} \right\rangle
        &= \sum_{k=1}^{n} \left\langle
        e_{ab}, e_{ik} \right\rangle
        \left\langle e_{cd}, e_{kj} \right\rangle \\
        &= \sum_{k=1}^{n}
        \delta_{ai}\delta_{bk}\delta_{ck}\delta_{dj} \\
        &= \delta_{ai}\delta_{bc}\delta_{dj}.
    \end{align*}
    The two expressions agree. Since the matrix unix form an orthonormal
    basis, this linearly extends to all matrices.
\end{proof}

\subsection{Crossed product quantum sets}

\textcite{tobolski26} have introduced the notion of crossed product quantum
sets. Suppose we are given a finite-dimensional \(C^{*}\)-algebra
\( \tilde{B}
\) along with a dual group action \( \hat{a}\colon \hat{G} \rightarrow
\operatorname{Aut}(\tilde{B}) \) where \( G \) is finite and
abelian. One can
then consider the crossed product
\[
    B \coloneqq \tilde{B} \rtimes_{\hat{\alpha}} \hat{G}.
\]

If \( \tilde{B} \) is equipped with a faithful functional \( \tilde{\psi} \)
in such a way as to make \( (\tilde{B}, \tilde{\psi}) \) a quantum
set, one can define a new faithful functional on \( B \) to make it a
quantum set, as described below.

\begin{definition}[Crossed product quantum set,
    \textcite[Definition~3.1]{tobolski26}]
    Let \( (\tilde{B}, \tilde{\psi}) \) be a quantum set and let \(
        \hat{\alpha}\colon \hat{G} \rightarrow \operatorname{Aut}(\tilde{B},
    \tilde{\psi}) \) be an action of the dual of a finite abelian
    group \( G \). The \emph{crossed product quantum set} is defined as
    \[
        (\tilde{B}, \tilde{\psi}) \rtimes_{\hat{\alpha}} \hat{G} \coloneqq
        \left( \tilde{B} \rtimes_{\hat{\alpha}} \hat{G},
        \psi_{\rtimes} \right),
    \]
    where \( \tilde{B} \rtimes_{\hat{\alpha}} \hat{G} \) is the
    crossed product \(C^{*}\)-algebra and \( \psi_{\rtimes} \) is the
    functional given by
    \[
        \psi_{\rtimes} \left( \sum_{\chi \in \hat{G}}
        b_{\chi}u_{\chi} \right) = n \tilde{\psi}(b_{1})
    \]
    with \( n = | \hat{G} | \).
\end{definition}

The proof that this definition is sensible, i.e. the newly defined
object is indeed a quantum set, can be seen in \textcite{tobolski26},
Proposition 3.4.

% TODO: Prove this lemma
\begin{lemma}
    Let \( \hat{\alpha}\colon \hat{\mathbb{Z}}_{n} \rightarrow
    \operatorname{A} \) be
    the action of cyclically shifting elements. Then
    \[
        (M_{n}, n \operatorname{Tr}) \cong (\mathbb{C}^{n} \rtimes_{\alpha}
        \hat{\mathbb{Z}}_{n}, \psi_{\rtimes}).
    \]
\end{lemma}

\section{Quantum graphs}

A simple graph is classically defined as a pair \( (V, E) \) with \( V \)
serving as the \emph{set of vertices} and \( E \subseteq [V]^{2} \)
serving as the \emph{set of edges}. Such a graph could be interpreted
as a matrix \( A = [a_{ij}] \in \{0, 1\}^{n \times n} \) with \(
\{v_{i}, v_{j}\} \in E \Leftrightarrow a_{ij} = 1 \); such a
representation is often called the \emph{adjacency matrix}. The key
property of such matrices is that they are invariant with respect to
entrywise multiplication. In turn we can use this property to
generalise graphs on sets to graphs on quantum sets.

The following definitions are due to \textcite{tobolski26}.
% TODO: Tell where this notion was first introduced
\begin{definition}[Quantum adjacency matrix,
    \textcite[Definition~2.5]{tobolski26}]
    Let \( (B, \psi) \) be a quantum set. A \emph{quantum adjacency
    matrix} on this quantum set is a linear map \( A\colon B \rightarrow B
    \) which is Schur-idempotent and *-preserving, by which we mean
    that it satisfies
    \[
        m(A \otimes A)m^{*} = A,
        \qquad
        A(b^{*}) = A(b)^{*}
        \quad \text{for all } b \in B.
    \]
\end{definition}

\begin{definition}[Quantum graph, \textcite[Definition~2.5]{tobolski26}]
    A \emph{quantum graph} on \( (B, \psi) \) is given by a quantum
    adjacency matrix on this quantum set.
\end{definition}

For all intents and purposes we will identify both notions and
simply call an
appropriate \( A\colon B \rightarrow B \) both a quantum adjancenty matrix and a
quantum graph.

\begin{remark}
    There are other equivalent definitions of a quantum graph, namely those
    based on
    \begin{itemize}
        \item an \emph{edge projection}, i.e. a projection \(
            \tilde{A} \in C(X) \otimes C(X)^{\operatorname{op}} \)
            \textcite{musto18},
        \item
    \end{itemize}
\end{remark}

\begin{definition}[\textcite[Definition~4.6{tobolski26}]]
    Let \( A \) be a quantum graph on a quantum set \( (B, \psi) \).
    We call the quantum graph
    \begin{itemize}
        \item \emph{loopfree} if \( m(A \otimes
            \operatorname{Id})m^{*} = 0, \)
        \item \emph{undirected} if \( A = A^{*} \), where the adjoint
            is taken with respect to the inner product induced by
            \( \psi \).
    \end{itemize}
    We say that a voltage quantum graph \( \tilde{A} = \left\{
    \tilde{A}_{g} \right\}_{g \in G} \) has these properties if the
    same equations hold for \( \sum_{g \in G} \tilde{A}_{g} \),
    respectively.
\end{definition}

\section{Isomorphisms and Quantum Isomorphisms of Quantum Sets}

\begin{definition}[Quantum isomorphisms,
    \textcite[Definition~2.8]{tobolski26}]
    Let \( (B_{1}, \psi_{1}) \) and \( (B_{2}, \psi_{2}) \) be
    quantum sets and
    let \( A_{1}\colon B_{1} \rightarrow B_{1} \) and \( A_{2}\colon B_{2}
    \rightarrow B_{2} \)
    be quantum graphs on \( (B_{1}, \psi_{1}) \), \( (B_{2},
    \psi_{2}) \), respectively.
    \begin{itemize}
        \item The two quantum sets are \emph{isomorphic} if there is
            a \( * \)-isomorphism
            \( \phi\colon B_{1} \rightarrow B_{2} \) such that \( \psi_{2}
            \phi = \psi_{1} \).
            We denote by \( \operatorname{Aut}(B, \psi) \) the group
            of quantum set
            automorphisms of a quantum set \( (B, \psi) \).
        \item The two quantum graphs are \emph{isomorphic} if there
            is an isomorpshism
            \( \phi\colon B_{1} \rightarrow B_{2} \) of the underlying
            quantum sets such that
            \( \phi A_{1} = A_{2} \phi \). We denote by \(
            \operatorname{Aut}(A) \) the
            group of quantum graph automorphisms of a quantum graph \( A \).
        \item The two quantum sets are \emph{quantum isomorphic}, written
            \( (B_{1}, \psi_{1}) \cong (B_{2}, \psi_{2}) \), if there is
            a finite-dimensional Hilbert space \( \mathcal{H} \) and a
            unital \( * \)-homomorphism
            \[
                \rho\colon B_{1} \rightarrow B_{2} \otimes B(\mathcal{H}),
            \]
            such that the map
            \[
                p\colon L^{2}(B_{1}) \otimes \mathcal{H} \rightarrow
                L^{2}(B_{2}) \otimes \mathcal{H}, \; a \otimes \xi
                \mapsto \rho(a)(1 \otimes \xi)
            \]
            is unitary as an operator between Hilbert spaces.
        \item The two quantum graphs are \emph{quantum isomorphic},
            written \( A_{1} \cong_{q} A_{2} \), if there is a
            quantum isomorphism \( \rho\colon B_{1} \rightarrow B_{2}
            \otimes B(\mathcal{H}) \) between the underlying quantum
            sets such that
            \[
                \rho A_{1} = (A_{2} \otimes
                \operatorname{Id}_{B(\mathcal{H})}) \rho
            \]
            holds.
    \end{itemize}
\end{definition}
